\cotefile{157-4}
\scanpage{228}
\mspage{102}
% archive sheet number: 227

Passons aux conditions supplémentaire\unsure{s} $e'$, $f$, $g$ exigées dans le cor. 7.4. La condition \unsure{$e)$|$e')$}
(\deleted{stabilité} \inserted{\unsure{comm.}} des $\underline{\lim}$ filtrantes au chgt de base) est satisfaite dans $A^\wedge$ et dans Cat. La
condition $f)$ dit que pour deux sous-objets $A_1,A_2$ \unsure{d'un} \deleted{$A$} objet $B$, si $A_0=A_1\cap A_2$,
alors

\[
 A_1\mathbin{\amalg}_{A_0}A_2\unsure{\longrightarrow|\text{---}}B
\]

est mono. Je doute que c'est vrai dans Cat, et ce n'est pas grave, car c'est \unsure{:|à} $A^\wedge$,
non \unsure{:|à} Cat, qu'on voudra appliquer 7.4. --- et c'est évidemment vrai dans $A^\wedge$, \deleted{\unsure{et A, est}} \inserted{et}
dans tout topos, l'étant dans (Ens). Enfin, la stabilité des monos par cochgt. de base
est sûrement fausse dans Cat, mais vraie dans $A^\wedge$ -- donc c'est OK.

Comme $W$ stable par $\underline{\lim}$ filtrantes
($W(\unsure{6|\sigma|o})$), et que $i_A$ y commute, $W_A$ est également stable. De $=$, \deleted{l'image inverse} \inserted{$W_A$}
\deleted{une partie} est $\unsure{\mathrm L|\mathrm U}$-accessible, \deleted{\unsure{par un}} en vertu du

\uline{Lemme} 8.1. Soient $\unsure{\mathcal{M}|M}\deleted{\unsure{N|M}},\mathcal{M}'$ deux catégories satisfaisant $a)\ b)\ c)\ d)\ e)$, et $\varphi:\mathcal{M}\longrightarrow \mathcal{M}'$
un foncteur commutant aux $\underline{\lim}$ filtrantes. Alors pour tte partie $\mathrm L$-accessible $\unsure{M'|N'}$ de $\mathcal{M}'$, $\varphi^{-1}(M')$ est $\mathrm L$-accessible.
\scanpage{229}
\mspage{103}
% archive sheet number: 228

De façon précise, \unsure{soit} $\pi$ un cardinal tel que
$\unsure{M'|N'}$ soit $\pi$-accessible et que $\varphi$ applique
$\unsure{\mathcal{M}|M}^\pi$ dans ${\mathcal{M}'}^\pi$ (\unsure{Par exemple|par exemple|Par Exemple}, pour cette dernière
condition, il suffit $\pi=2^c$ pour $c$ assez
grand, cf loc cit. 9.14 p. 153) Alors \deleted{$\varphi^{-1}(\unsure{N|M}$}
$M=\varphi^{-1}(M')$ est $\pi$-accessible.

Vérification immédiate.

On a donc le

\uline{Lemme} 8.2. Soit $A$ dans Cat. Alors $(A^\wedge,W_A)$
satisfait aux conditions préliminaires du
cor. \unsure{7|1}.4., ainsi qu'aux conditions suivantes

\[
 \deleted{\unsure{\mathcal C\subset|\mathrm{Le}}}
 \underset{\substack{\shortparallel\\\text{\unsure{ens.} des}\\\text{mono}\\\text{de } A^\wedge}}
 {\mathrm{Mon}(A^\wedge)}\subset\mathrm{Cof}_{W_A}
 \tag*{(8.2.1)}
\]

\[
 \deleted{\unsure{\mathrm i\in|\mathrm{ces}}}\quad \mathrm{Mon}\cap W_A
 \quad\text{stable\unsure{s} par coch. de base}
 \tag*{\unsure{(8.2.2)|(8.1.b)}}
\]

en d'autres termes (compte tenu que Mon est stable par coch. de base) pour
un carré cocartésien

\diagram{\begin{tikzcd}[ampersand replacement=\&,column sep=small]
 {(8.2.3)} \& A\arrow[rr,hook,"i"]\arrow[dr,"\alpha"'] \& \& B\arrow[dr,"\beta"] \& {}\\
 {} \& {} \& A'\arrow[rr,hook,"i'"] \& {} \& \deleted{\unsure{B|A}}B'
\end{tikzcd}}

avec $i$ mono, si $\alpha\in W_A$ (resp. $i\in W_A$),
on a $\beta\in W_A$ (resp. $i'\in W_A$). En fait
on a mieux : \deleted{le \unsure{dit|carré} carré est $W$-cocartésien}

\[
 \begin{aligned}
 &\deleted{i\in W_A\unsure{\Longleftrightarrow|\Longrightarrow}}\\
 &i\in W_A\unsure{\Longleftrightarrow|\Longrightarrow} i'\in W_A,\qquad
 \alpha\in W_A\Longleftrightarrow\beta\in W_A.
 \end{aligned}
 \tag*{(8.2.4)}
\]
\scanpage{230}
\mspage{104}
% archive sheet number: 229

Il reste à vérifier (8.24.), qui a été vu dans $\overline{\underline{\unsure{\mathrm{XII}|\mathrm{XV}|\mathrm{XIV}}}}$\ldots\quad / On conclut donc de 7.4. $1^\circ,\unsure{1|2}^\circ,3^\circ$

\uline{Corollaire} 8.3. Posons \unsure{\S|$\xi$}

\[
 \left\{\begin{aligned}
 C_{\deleted{\unsure{A}}}&=\mathrm{Mon}(A^\wedge),&
 \mathrm{TC}_{\deleted{\unsure{A}}}&=C\cap W_{\unsure{A|a}}\\
 F&=(\mathrm{TC})_\ast,&\mathrm{TF}&=C_\ast.
 \end{aligned}\right.
 \tag*{(8.3.1\unsure{.})}
\]

Alors $(C,\mathrm{TF})$ et $(\mathrm{TC},F)$ sont des couples
de Quillen, à factorisation\unsure{s} fonctorielle\unsure{s} (NB
$C$ et $\mathrm{TC}$ \deleted{\unsure{étant}} \inserted{de plus} Q-engendrées \unsure{à p.|i.e.} par des
petites parties de $\mathrm{Fl}\,\widehat A$). Pour que
$(W_A,C,F)$ soit un triple de Quillen \unsure{clos},
il \unsure{f.} et s. qu'on ait

\[
 \underset{\substack{\shortparallel\,\mathrm{def}\\(\mathrm{Mon})_\ast}}{\mathrm{TF}}\subset W_A\ ,
 \tag*{(8.3.2)}
\]

\deleted{\unsure{en}} i.e. \deleted{\unsure{c'est à|on ait}} que toute flèche $f:X\to Y$ \inserted{de type injectif}
de $A^\wedge$, \inserted{(i.e. qui fait de $X$ un objet injectif de $A^\wedge/Y$)} soit dans $W_A$.
Quand cette condition est satisfaite, le triple de Quillen $(W_A,C,F)$ est propre à gauche, par 8.2.1.
Il n'est pas clair qu'elle soit forcément propre -- totalement, i.e. qu'on ait

\[
 F\subset\mathrm{Fib}_{W}\quad (?)\ .
 \tag*{(8.3.3)}
\]

Nous y reviendrons par la suite. Pour l'instant, la chose essentielle, c'est d'examiner la condition (8.3.2) sur $A$.
\scanpage{231}
\mspage{105}
% archive sheet number: 230

\uline{Théorème 8.4.} Soit $W$ un localiseur \inserted{fond.} dans Cat
satisfaisant à $W(1,2,3,\deleted{\unsure{5|4}},6,7,8)$ \inserted{notations de $\overline{\underline{\mathrm{XII}}}$} (i.e. tous les
axiomes envisagés jusqu'à présent, sauf \inserted{peut-être} celui
(4) des fibrations),
\inserted{Supposons de plus $W$ stable par fact. directs (\unsure{renforcement} $W(1\,\mathrm{bis})$ de l'axiome de saturation \unsure{modérée|modéré} $W(1)$}
et soit $A$ un objet de Cat, et $W_A=\unsure{v_A|i_A}^{-1}(W)$ le localiseur
de $A^\wedge$ déduit de $W$. Avec les notations de 8.3,
\deleted{pour qu'on ait $\mathrm{TF}$ la condition suivante}

% slanted
\margin{\uline{NB} Pour la validité de $1^\circ),\unsure{3|2}^\circ)$ \deleted{\unsure{et le cor.}}, on n'a pas besoin de $W(6)$ ($\underrightarrow{\lim}$), ni de $W(8)$ (accessibilité), ni même de $W(7)$ (axiome \unsure{du|des} carré\unsure{s} cocartésien\unsure{s}). Ce dernier est \unsure{utile|mis à contribution} pour assurer \deleted{que} le $a)$ du cor. 8.5. Mais pour la partie $2^\circ)$ du Th. 8.4., on a besoin de \unsure{$g)$|$3^\circ$} \uline{tous} et \unsure{seulement} les axiomes du \unsure{localiseur} fondamental $W$ \unsure{à l'exc. de|à} $W(4)$.}

\textcircled{$1^\circ$}) \uline{Les conditions} \deleted{suivantes} sont équivalentes

$a)$ On a (8.3.2)

\[
 \mathrm{TF}\subset W_A\ ,
\]

i.e. toute flèche de type injectif dans $A^\wedge$ est dans $W_A$

$b)$ Soit

\[
 L=L_A \tag*{(8.4.1)}
\]

l'élément de Lawvere dans $A^\wedge$, donc le préfaisceau

\[
 \begin{aligned}
 a\longmapsto{}&\text{ensemble des sous-préfaisceaux}\\
 &\text{de } a\ \text{(i.e. ensemble}\\
 &\text{des \deleted{cribles} ouverts de}\\
 & A_{/a},\ \text{i.e. des cribles}\\
 &\text{de } A_{/a})\ ,
 \end{aligned}
\]

et considérons la flèche $L\longrightarrow e_A$ (objet final de $A^\wedge$). On veut que cette flèche soit \deleted{dans}
$W_A$-asphérique, i.e. soit universellement dans $W_A$ -- ou ce qui revient au $=$, que pour tt $a\in A$, $L_A\times a$
$(=L_{A_{/a}})$ soit un objet $W_A$-asphérique
\scanpage{232}
\mspage{106}
% archive sheet number: 231

de $A^\wedge$ (ou, ce qui revient au $=$, \deleted{que} de $(A_{/a})^\wedge$).

\inserted{$c)$ Il existe un objet $I$ de $A^\wedge$ tel que $I\to e_A$ soit
$W_A$-asphérique, \deleted{et que} avec deux sections $\delta_0,\delta_1$ de $I$ sur
$e_A$, \unsure{avec|telles} $\mathrm{Ker}(\delta_0,\delta_1)=\varnothing_A$ (préfaisceau vide).}

$d)$ Les catégories localisées $A_{/a}=B$ sont des
\uline{catégories-test} pour $W$ (cf Pursuing Stacks, \S\,65, notamment pages 172--173)
ce qui revient à dire que l'on a ceci pour chacune des \deleted{\unsure{ca}} catégories $B=A_{/a}$ :

% slanted
\margin{\uline{NB} La condition $d)$, d'apparence très forte, est équivalente à la condition d'\unsure{apparence|allure} plus \unsure{anodine} $d')$ Le foncteur $j_A:\mathrm{Cat}\to A^\wedge$ transforme \unsure{flèches|éléments} à \unsure{fct adjoint|fonc. adjoints} en flèches $W_A$-asphériques. -- et il suffit $=$ de l'exiger pour $C\to e$ dans Cat, où $C$ \deleted{est} un objet de Cat ayant objet final.}

\deleted{(8.4.2)} Les foncteurs

\[
 \left\{\begin{aligned}
 i_B &: B^\wedge\longrightarrow\mathrm{Cat}
 & (F\unsure{\mapstochar\relbar|\longmapsto} B_{/F})\\
 j_B &: \mathrm{Cat}\unsure{\mathrel{\text{---}}|\longrightarrow} B^\wedge
 &\deleted{(C\longmapsto\mathrm{Hom}_{\mathrm{Cat}}(}\\
 &&(C\mapsto(b\longmapsto\mathrm{Hom}(B_{/b},C)))
 \end{aligned}\right.
 \tag*{(8.4.2)}
\]

\deleted{qui} où $j_B$ est l'adjoint à droite de $i_B$) sont tous les deux ``model preserving'' i.e.

\[
 \left\{\begin{aligned}
 W_B&=i_B^{-1}(W) &&\text{(par définition de }W_B\text{ en fait)}\\
 W&=j_B^{-1}(W_B)
 \end{aligned}\right.
 \tag*{(8.4.3.)}
\]

et les foncteurs \deleted{local} localisés

\diagram{\begin{tikzcd}[ampersand replacement=\&]
 \mathrm{Hot}_{B,W}\arrow[r,shift left,"\overline{i_B}"] \&
 \mathrm{Hot}_{W}\arrow[l,shift left,"\overline{j_B}"]
\end{tikzcd}}

sont des équivalences de catégories. De plus, ces équivalences de catégories sont
quasi-inverses l'une de l'autre, les morphismes d'adjonction étant déduits
des morphismes d'adjonction pour les couples foncteurs \inserted{adj.} $(i_B,j_B)$, ces morphismes d'adjonction
étant des \unsure{flèches|élts} dans $W$ respt. dans $W_A$.
\scanpage{233}
\mspage{107}
% archive sheet number: 232

$2^\circ)$ Si les conditions équivalentes du $1^\circ)$ sont satisfaites, alors le triple

\[
 (W_A,\ C=\mathrm{Mon}_{A^\wedge},\ F=(C\cap W_A)_\ast)
\]

est un triple de Quillen \inserted{(clos)} propre à gauche

$3^\circ)$ Si les conditions \inserted{équiv.} du $1^\circ)$ sont satisfaites,
et si de plus $A$ est $W$-asphérique, alors \inserted{(et alors seulement)}
$A$ elle-même est une \uline{catégorie-test} \uline{pour}
$W$, ce qui revient à dire que les \deleted{\unsure{propri}}
conditions sur la cat. $B$ dans $d)$ sont satisfaites quand $B$ est \uline{soit} $A$, soit
une \deleted{\unsure{des|ses}} catégorie localisée $A_{/a}$.

\uline{Corollaire} \inserted{8.5} : Les conditions équivalentes du $1^\circ)$\inserted{, et du $3^\circ)$,} sont satisfaites dans chacun des cas
suivants (et ceci quel que soit $W$ satisfaisant aux hypothèses \unsure{dites|étant})

$a)$ $A=\mathord{\tikz[baseline=-.05ex,x=1em,y=1em]{\draw[line width=.055em] (0,0)--(.83,.03)--(.52,.67)--cycle; \draw[line width=.05em] (.13,.03)--(.45,.57);}},\mathord{\tikz[baseline=-.05ex,x=1em,y=1em]{\draw[line width=.055em] (0,0)--(.83,.03)--(.52,.67)--cycle; \draw[line width=.05em] (.13,.03)--(.45,.57);}}^{\sim},\mathord{\tikz[baseline=-.05ex,x=1em,y=1em]{\draw[line width=.055em] (.08,0)--(.78,.02)--(.69,.61)--(.02,.56)--cycle; \draw[line width=.05em] (.17,0)--(.11,.56);}},\mathord{\tikz[baseline=-.05ex,x=1em,y=1em]{\draw[line width=.055em] (.08,0)--(.78,.02)--(.69,.61)--(.02,.56)--cycle; \draw[line width=.05em] (.17,0)--(.11,.56);}}^{\sim}$ (catégories des
simplexes \inserted{combinatoires} \unsure{standart|standard}, ordonnés ou non,
et des cubes combinatoires standart, ordonnés ou non \ldots)

% slanted
\margin{$\mathord{\tikz[baseline=-.05ex,x=1em,y=1em]{\draw[line width=.055em] (.08,0)--(.78,.02)--(.69,.61)--(.02,.56)--cycle; \draw[line width=.05em] (.17,0)--(.11,.56);}}^{\sim}$ est \unsure{essentiellement} la cat. $A$ ``universelle'' stable par produits \inserted{(finis)} et contenant un ``intervalle séparant'' $(\unsure{\overline I|I},\delta_0,\delta_1)$ comme dans $b)$}

$b)$ $A$ est stable par produits finis,$^{\ast}$ et contient un objet $I$ tel que $I$
ait deux sections $\delta_0,\delta_1$ sur l'objet final, avec $\delta_0,\delta_1$ disjointes i.e.
$\mathrm{Ker}(\delta_0,\delta_1)=\varnothing_{A^\wedge}$ (sous-\deleted{préfaisceau} vide du préfaisceau final).

% slanted
\margin{en fait, il suffit que $I\to e$ soit \unsure{quarrable|carrable} i.e que \deleted{\unsure{$e$}} $A$ \deleted{\unsure{stable}} ait un objet final et que les produits $I\times a$ ($a\in\mathrm{Ob}\, A$) existent dans $A$.}
\scanpage{234}
\mspage{108}
% archive sheet number: 233

Le cas $b$ du corollaire est justiciable de la condition $c)$ du $1^\circ)$. (De plus $A$ ayant un objet final, est $W$-asph., donc on est dans les conditions du $3^\circ)$
Les cas $\deleted{\unsure{\mathord{\tikz[baseline=-.05ex,x=1em,y=1em]{\draw[line width=.055em] (.08,0)--(.78,.02)--(.69,.61)--(.02,.56)--cycle; \draw[line width=.05em] (.17,0)--(.11,.56);}}^\sim}}\mathord{\tikz[baseline=-.05ex,x=1em,y=1em]{\draw[line width=.055em] (0,0)--(.83,.03)--(.52,.67)--cycle; \draw[line width=.05em] (.13,.03)--(.45,.57);}},\mathord{\tikz[baseline=-.05ex,x=1em,y=1em]{\draw[line width=.055em] (0,0)--(.83,.03)--(.52,.67)--cycle; \draw[line width=.05em] (.13,.03)--(.45,.57);}}^{\sim},\mathord{\tikz[baseline=-.05ex,x=1em,y=1em]{\draw[line width=.055em] (.08,0)--(.78,.02)--(.69,.61)--(.02,.56)--cycle; \draw[line width=.05em] (.17,0)--(.11,.56);}}$ et $\mathord{\tikz[baseline=-.05ex,x=1em,y=1em]{\draw[line width=.055em] (.08,0)--(.78,.02)--(.69,.61)--(.02,.56)--cycle; \draw[line width=.05em] (.17,0)--(.11,.56);}}^\sim$, qui tombent en fait dans le cas $b)$ sont traités dans P.S. \S\,34 (p. 47--50)

\uline{Démonstration du th. 8.4.} La partie essentielle, techniquement, est bien sûr $1^\circ)$.$^\ast$ On a les implications $\pm$ tautologiques

% slanted
\margin{La partie $2^\circ$ en résulte, grâce à cor. 8.3.}

\[
 a\Longrightarrow b\Longrightarrow c\quad\deleted{\text{\unsure{immédiat|immédiate}}},
\]

où $a\Longrightarrow b$ provient du fait que $L$ est évidemment un objet injectif de $A^\wedge$. (L'objet de Lawvere d'un topos quelconque est un objet injectif). Il reste à établir

\[
 c\Longrightarrow d\Longrightarrow a.
\]

Que l'on ait $c\Longrightarrow d$ (et $c\Longrightarrow b$) est établi dans PS, voir le résumé des \unsure{critères|catég.} $W$-test de PS p. \unsure{168|468}--179\unsure{$'$} ; l'énoncé pertinent ici est prop. 3 page 173.

Il reste à établir $d)\Longrightarrow a)$ -- mais à présent on joue sur du velours !
Je \deleted{\illegible{1}} note ici le lemme suivant \deleted{\unsure{dont|qu'est}} la place serait plutôt à la fin du \S\,4 \unsure{des} Généralités sur les triples de Quillen :
\scanpage{235}
\mspage{108}
% archive sheet number: 234

\deleted{Notons d'ailleurs que pour prouver que \deleted{\unsure{toute}} $f\in\mathrm{TF}$ est dans $W_A$, et $=$ est
$W_A$-asphérique, on est ramené aussitôt au cas où son but est un objet $a$ de $A$,
\diagram{\begin{tikzcd}[ampersand replacement=\&]
 \deleted{\unsure{Y|F|X}}X\arrow[d] \& X_a\arrow[l]\arrow[d]\\
 Y \& a\arrow[l]
\end{tikzcd}}
donc l'objet final de la catégorie $A_{/a}$. On est ainsi ramené au cas d'une
catégorie $\deleted{\unsure{B=A_{/a}}}$ $B=\unsure{A_{/a}|A_a}$ avec objet final,
vu que l'hyp. $d)$ est stable par passage de $A$ à une localisée $A_{/a}$.

À prouver que dans le cas où $A$ a un objet final, donc est $W$-asphérique,
pour tout $X$ dans $A^\wedge$, $X\xrightarrow{f}e$ se factorise en
$X\hookrightarrow X'\to e$, avec $X\to X'$ mono et $X'\longrightarrow e$ dans $W_A$.
Je vois}

\unsure{Réellement|Brutalement|Décidément}, mon projet de démonstration se \unsure{modifie|modifia}, car bien que
$d)$ ait l'air très, très fort, comme condition,
\deleted{\illegible{1}} telle quelle je ne vois pas, \unsure{pratiquement|finalement},
comment en déduire $a)$. Mais en fait
\deleted{\unsure{alors}} je \unsure{peux|puis|vais} considérer que l'équivalence
de $b,c,d)$ est mise là pour mémoire, elle est établie dans loc cit,
donc je n'ai à me préoccuper que de l'implication $b\Longrightarrow a$, qui va résulter
\scanpage{236}
\mspage{109}
% archive sheet number: 235

\uline{Lemme} 8.6. Soit $\unsure{\mathcal{M}|M}$ une catégorie, $(C,\mathrm{TF})$
un précouple de Quillen dans $\mathcal{M}$, et $W\subset\mathrm{Fl}\,\mathcal{M}$ une partie de $\mathrm{Fl}\,\mathcal{M}$, stable par facteurs directs.
\deleted{Pour qu'on ait $\mathrm{TF}\subset W$, il faut et \unsure{s.} que toute flèche $f$ \unsure{a}} \inserted{Conditions équivalentes}

$a)$ \deleted{\unsure{tt}} $\mathrm{TF}\subset W$

$b)$ toute flèche $f$ dans $\mathcal{M}$ se factorise en $f=pi$, avec $i\in C$, $p\in W$.

$b')$ Comme $b)$, avec $f$ dans $\mathrm{TF}$.

Démonstration : les implications $a)\Rightarrow b)\Rightarrow b')$
sont tautologiques, et ne font pas appel à l'hypothèse \inserted{de stabilité} sur $\unsure{W_a|W.|W_o}$
Mais pour un $f\in\mathrm{TF}$ donné, pour voir qu'il est dans
$W$, \deleted{on utilise l'hyp} si $W$ est stable par facteurs directs, il suffit en effet de
le factoriser en $pi$, avec $i\in\unsure{\mathrm{TC}|C}$, $p\in W$.
En effet, \deleted{\unsure{on}} voit aussitôt qu'alors
$f$ est facteur direct de $p$ (car $i\in C$) et comme $p\in W$, on a $i\in W$, \unsure{qed.}

Revenons à 8.4. dém de $d)\Rightarrow a)$.

\deleted{On veut prouver que $\mathrm{TF}\subset W$, et on est ramené à prouver que toute flèche $f$ de $A^\wedge$ se factorise en $pi$, avec $i\in C$ i.e. $i$ mono, et $p\in W_A$}
\scanpage{237}
\mspage{111}
% archive sheet number: 236

du lemme 8.6., qui nous ramène à prouver que toute flèche $f:X\unsure{\mathrel{\text{---}}|\longrightarrow}Y$ dans $A^\wedge$ se factorise en $pi$, avec $i$ mono et $p\in W_A$. On fait l'argument du mapping-cône, en partant de

\diagram{\begin{tikzcd}[ampersand replacement=\&,column sep=small]
e_A\amalg e_A \arrow[rr,hook,"{(\delta_0,\delta_1)}","{\mathrm{mono}}"']\arrow[dr,dash] \& \& \deleted{\unsure{A}}I\arrow[dl,"{(W_A)_{\mathrm{univ}}}"]\\
\& e_A \&
\end{tikzcd}}

d'où pour tout $X$ dans $\widehat A$, en prenant le produit par $X$

\diagram{\begin{tikzcd}[ampersand replacement=\&,column sep=small]
X\amalg X\arrow[rr,hook,"{(\deleted{\unsure{\delta}}\delta_0^X,\deleted{\unsure{1}}\delta_1^X)}"]\arrow[dr,"{(\mathrm{id},\mathrm{id})}"] \& \& \deleted{\unsure{I}}X\times\deleted{\unsure{X}}I\arrow[dl,dash,"{\mathrm{pr}_1^X\in W_A}"]\\
\& X \&
\end{tikzcd}}

et on construit le mapping-\unsure{cone|cône} de $f$ par le diagramme cocartésien

\diagram{\begin{tikzcd}[ampersand replacement=\&,column sep=large,row sep=large]
\& X\deleted{\unsure{I}}\arrow[r,"f"]\arrow[dr,phantom,"{\scriptstyle\mathrm{cocart}}" description]\arrow[d,"{\delta_0^X}"'] \& \deleted{\unsure{Z}}Y\deleted{\unsure{(f)|f}}\arrow[d,"\beta"] \& {\substack{p_f\ \text{déduit}\\ \text{de }\mathrm{pr}_1^X:X\times I\to I}} \\
{(8.6.1)}\quad X\arrow[r,"{\delta_1^X}"]\arrow[rr,bend right=24,dashed,"{i_f}"'] \& X\times I\arrow[u,bend right=28,dashed,"{\mathrm{pr}_1^X}"']\arrow[r,"{f'}"] \& Z(f)\arrow[u,bend right=36,dashed,"{p_f}"'] \& {,}
\end{tikzcd}}

la factorisation cherchée de $f$ est

\[
f=p_f i_f\qquad\text{où }\ i_f=f'\delta_1^X .
\]

Ici, comme $\mathrm{pr}_1^X\in W_A$, on en conclut que
\scanpage{238}
\mspage{112}
% archive sheet number: 237

$\delta_0^X,\delta_1^X$ (qui en sont des sections) sont dans $W_A$, comme $\delta_0^X$ est mono on a donc $\delta_0^X\in\mathrm{TC}$, donc ($\mathrm{TC}$ étant stable par \unsure{coch.\textsuperscript{t}|coch\textsuperscript{t}|coch\textsuperscript{gt}} de base) $\beta\in W_A$, donc enfin $p_f$ (qui est une rétraction de $Z(f)$ sur $Y$) est aussi dans $W_A$.

Il reste à voir que $i_f$ est mono, ce qui se voit sur le diagramme plus complet (cf $\overline{\underline{\mathrm{XII}}}$ p. 230\deleted{\unsure{)}}) que j'ai sous les yeux : $i_f$ apparaît comme composée de deux mono. $X\longrightarrow Y\amalg X$ et $Y\amalg X\longrightarrow Z(f)$

\inserted{\unsure{Le Th.|le th.} 8.4. est démontré.}

\unsure{.....|....} Je vais quand même isoler l'argument formel \unsure{supplémentaire} (que j'ai utilisé surtout sous forme duale) dans un

\uline{Lemme} 8.7. Soit $\unsure{\mathcal{M}|M}$ une catégorie, $W$ un localiseur modérément saturé dans $\mathcal{M}$, $C$ une partie stable par composition et \deleted{\unsure{contenant les}} \inserted{par cochangt.} de base,

On se propose, pour $X$ fixé dans $\mathcal{M}$, de trouver une factorisation fonctorielle de \unsure{tte|toute} flèche $f:X\unsure{\mathrel{\text{---}}|\longrightarrow}Y$ (i.e. du morphisme structural d'un $Y$ dans $X\backslash\mathcal{M}$) en $p_f i_f$, avec $i_f\in C$, $p_f\in W$. On suppose \deleted{que} pour ceci qu'on a
\scanpage{239}
\mspage{113}
% archive sheet number: 238

\unsure{-|.} une \deleted{telle} factorisation de ce type pour $X\amalg X\xrightarrow{(\mathrm{id},\mathrm{id})}X$, i.e. un objet $I_X$ et une \deleted{\unsure{diagr.}} factorisation

\diagram{\begin{tikzcd}[ampersand replacement=\&,column sep=large]
X\amalg X\arrow[r,hook,"{i_X}","{\substack{\Vert\\(\delta_0^X,\delta_1^X)}}"']\arrow[rr,bend right=24,dashed,"{(\mathrm{id}_X,\mathrm{id}_X)}"'] \& I_X\arrow[r,"{p_X}"] \& X,
\end{tikzcd}}

avec $i_X\in C$, $p_X\in W$. On \deleted{\unsure{f|y}} associe alors à tout $f$ \deleted{\unsure{son|un}} \unsure{mapping cone|mapping-cone} $C(f)$ comme dans le diagramme (8.6.1) (où il faut lire $I_X$ ald $X\times I$), et on a alors $p_f\in W$, $i_f\in C$ comme voulu. Mais il faut supposer, pour savoir que $i_f\in C$, que \deleted{\unsure{NB|N.B.} \unsure{les diagrammes} \unsure{co--}} $X\in \unsure{\mathcal{M}|M}_C$ (objets $C$-cofibrants) i.e que pour \unsure{$\forall\,Y$|tt $Y$}, $X\unsure{\rightarrowtail|\longrightarrow}\deleted{\unsure{X}}Y\amalg X$ est dans $C$. (De plus, on a supposé tacitement $\mathcal{M}$ stable par $\unsure{\varinjlim|\underline{\lim}}$ finies.)

Le diagramme complet est

\diagram{\[
\left\{\begin{gathered}
\begin{tikzcd}[ampersand replacement=\&,column sep=large,row sep=large]
\& X\arrow[r,"f"]\arrow[d,"{\mathrm{inj}_0^X}"]\arrow[dr,phantom,"{\scriptstyle\mathrm{cocart}}" description] \& Y\arrow[d]\\
X\arrow[r,"{\mathrm{inj}_1^X}"]\arrow[dr,dashed,"{\delta_1^X}"'] \& X\amalg X\arrow[d,"{i_X}"]\arrow[r,"{f''}"]\arrow[dr,phantom,"{\scriptstyle\mathrm{cocart}}" description] \& Y\amalg X\arrow[d,"{\unsure{j|J}}"]\\
\& I_X\arrow[uu,bend right=20,dashed,"{p_X}"']\arrow[r,"{f'}"'] \& Z(f)\arrow[uu,bend right=38,dashed,"{p_f}"']
\end{tikzcd}\\
f=p_f i_f\ ,\ \text{où }\ i_f=f'\circ\delta_1^X=j\circ(\underbrace{f''\circ\mathrm{inj}_1^X}_{\lambda}) .
\end{gathered}\right.
\tag*{(8.\unsure{7}.1)}
\]}
\scanpage{240}
\mspage{114}
% archive sheet number: 239

On a $i_X\in C$ par hypothèse, donc $j\in C$ (stabilité de $C$ par coch. de base), on a $\lambda:X\longrightarrow Y\amalg X$ dans $C$ car $X$ dans $\unsure{\mathcal{M}|M}_C$, donc le composé $i_f=j\circ\lambda$ est dans $C$ (stabilité de $C$ par composition) --- on gagne !

% slanted
\margin{\uline{Commentaires}}

Pour l'équivalence de $a)$ $b)$ $c)$ dans le th. 8.4, $1^\circ$, on n'a \deleted{\unsure{plus}} pas eu à faire appel aux \deleted{conditions} développements de PS, qui servent \deleted{à \unsure{établir}} seulement à établir l'équivalence des trois conditions en question avec la condition ``\unsure{modélisante|anodinisante}'' $d)$. Pour résumer :

$1^\circ)$ les implications

\[
a\Rightarrow(b\Leftrightarrow c\Leftrightarrow d)
\]

sont valables sous les seules conditions

\[
W(1,2\mathrm{bis},3)
\]

et \deleted{\unsure{l'équi--}} pour l'équivalence de $a$ aux autres conditions, il suffit de

\[
W(1\mathrm{bis},2\mathrm{bis},3)
\]

(où $W(1,\mathrm{bis})$ renforce $W(1)$ par la condition des fact. directs, tandis que $W(2\mathrm{bis})$ affaiblit $W(2)$, en exigeant seulement que les morphismes $W$-asph. soient dans $\unsure{W^{\circ}|W}$)

\deleted{\unsure{$2^\circ)$}} Les conclusions de 8.4. $3^\circ)$, relatives au cas où $A$ est $W$-asphérique, sont valables également sous les conditions $W(1,2\mathrm{bis},3)$
\inserted{\deleted{\unsure{sous|sans}} plus}
\scanpage{241}
\mspage{115}
% archive sheet number: 240

$2^\circ)$ Donc le corollaire 8.5, \deleted{\unsure{la partie|la p.}} \unsure{relativement} aux conditions $b)$ $c)$ $d)$ du $1^\circ$ (à l'exclusion du $a)$, qui suppose $W$ stable par fact. directs) est valable pour un $W$ satisfaisant

\[
W(1,2\mathrm{bis},3,\underset{\inserted{\substack{\text{carré}\\\text{cocartésien}}}}{7\mathrm{bis}})
\]

$3^\circ)$ Pour \deleted{\unsure{la|la c.}} validité de 8.4. $2^\circ$, i.e. pour avoir un triple de Quillen $(W_A,C,F)$, il nous faut \inserted{(presque)} tout l'éventail des conditions sur $W$

\[
W(1\mathrm{bis},\underline{2\mathrm{bis}},3,\deleted{5},\underset{\inserted{\text{limites}}}{6},
\underset{\inserted{\substack{\text{carré}\\\text{coc.}}}}{7\mathrm{bis}},
\underset{\inserted{\text{accessibilité}}}{8})
\]

% slanted
\margin{\uline{NB} On devrait \unsure{renuméroter} les axiomes $W(\ldots.)$,}

à l'exclusion de $W(4)$, et en se \unsure{contentant}\deleted{\unsure{ant}} pour $W(2)$ de la forme affaiblie $W(2\mathrm{bis})$ sans plus. (Même l'\unsure{axiome} de connexité $W(5)$ n'a pas servi, mais \deleted{\unsure{que}} \inserted{\unsure{l'avais|au vrai|dans l'état} \unsure{d'état|actuel}} des choses, \deleted{\unsure{peut-être que}} car il y a peu de chances qu'on aura jamais à travailler avec des $W$ qui n'y satisfont pas ! \deleted{\unsure{)}} Il se pourrait d'ailleurs, en présence des autres axiomes, ou seulement de $W(1\mathrm{bis},2\mathrm{bis},3,7\mathrm{bis})$, que si $W(5)$ n'est \uline{pas} satisfait, on ait $W=\mathrm{Fl}\,\mathrm{Cat}$ tout entier, donc $W_A=\mathrm{Fl}\,\widehat A$ !)
\scanpage{242}
\mspage{116}
% archive sheet number: 241

% slanted
\margin{cf \uline{$\overline{\mathrm{XVII}}$} version définitive (?)}

9. \uline{Révision des \unsure{positions|notations}, pour les axiomes d'un localiseur fondamental}

\uline{Loc} \unsure{(1)} (Saturation) \quad $a)$ Les \deleted{\unsure{identités}} isomorphismes sont dans $W$

$b)$ Si $f,g$ composables, et deux parmi $f,g,gf$ dans $W$, le $3^\mathrm{e}$ aussi.

$c)$ Si
\diagram{\begin{tikzcd}[ampersand replacement=\&,column sep=large]
X\arrow[r,phantom,shift left=1.7ex,"{\overset{f}{\unsure{\longrightarrow|\mathrel{\text{---}}}}}" description]
\arrow[r,phantom,shift right=1.7ex,"{\underset{g}{\unsure{\mathrel{\text{---}}|\longleftarrow|\longrightarrow}}}" description] \& Y
\end{tikzcd}},
si $gf$ et $fg$ sont dans $W$, alors $f$ et $g$ aussi.

\unsure{variante} renforcée

\uline{Loc} \unsure{(1bis)} : Loc (1) $+$ \quad $d)$ $W$ stable par facteurs directs

\uline{Loc} \unsure{(2)} (Axiome de l'objet final) : Si $X$ dans Cat a un objet final, alors $X$ est $W$-asph. i.e. $X\longrightarrow e$ est dans $W$.

\uline{Loc} \unsure{(3)} (Axiome des \deleted{foncteurs asphériques} \inserted{\unsure{localisation|localisations}, forme faible}) Soit $f:X\unsure{\mathrel{\text{---}}|\longrightarrow}Y$ tel que pour tout $y\in Y$, $X_{/y}$ soit $W$-asphérique, (i.e le morphisme $f_{/y}:X_{/y}\longrightarrow Y_{/y}$ déduit de $f$ par chgt. de base $Y_{/y}\longrightarrow Y$ est dans $W$). Alors $f\in W$.

Variante renforcée de \uline{Loc} (3)

\uline{Loc} \unsure{(3bis)} (Axiome de localisation, forme forte) : Soit $f:X\longrightarrow Y$ un morphisme de $\mathrm{Cat}_{/S}$, \deleted{\unsure{que|quel}} \deleted{\unsure{que}} $\forall\,s\in S$, $f_{/s}:X_{/s}\unsure{\mathrel{\text{---}}|\longrightarrow}Y_{/s}$ (déduit de $f$ par chgt de base $S_{/s}\longrightarrow S$) soit $\in W$. Alors $f\in W$.

[NB. Sauf erreur, il suffit de l'exiger si $X,Y$ sont \deleted{\unsure{des}} fibrés scindés sur $S$, et $f$ \inserted{\unsure{cartésien|scindien}} compatible avec les scindages --- en supposant Loc (1,2)

\deleted{\uline{Loc} (4) \unsure{-} (Axiome des $W$-fibrations) Variantes $a)$ $b)$ $c)$ \unsure{ou|on}, cf \uline{XIV}}
